% บทคัดย่อไทยและอังกฤษ อยู่บนสุดของคอลัมน์ซ้าย ความยาว 100-150 คำ
% หัวข้อกึ่งกลาง ตัวหนา ไม่มีเลขกำกับ

\noindent{\centering\bfseries บทคัดย่อ\par}

งานวิจัยนี้ศึกษาความสัมพันธ์ระหว่างเวลาที่นักศึกษาใช้โซเชียลมีเดีย ประเภทเนื้อหาที่เสพ
และการติดตามข่าวเชิงลบ กับภาวะซึมเศร้า ความวิตกกังวล และความเครียด
ซึ่งวัดด้วยแบบวัด DASS-21 ฉบับภาษาไทย
ข้อมูลเก็บด้วยแบบสอบถามออนไลน์จากนักศึกษาระดับปริญญาตรี มหาวิทยาลัยขอนแก่น 121 คน
และใช้วิเคราะห์ 107 คนที่ผ่านข้อดักความใส่ใจ
ผลการวิเคราะห์ด้วยสหสัมพันธ์เพียร์สันและการถดถอยพหุคูณพบว่า
ชั่วโมงใช้โซเชียลมีเดียต่อวันไม่สัมพันธ์กับมิติใดอย่างมีนัยสำคัญ
(\stat{r} อยู่ระหว่าง −0.15 ถึง −0.09)
และเมื่อควบคุมการนอน ความกดดันจากการเรียน ผลการเรียน และสภาพการเงินแล้ว
ตัวแปรโซเชียลมีเดียเพิ่มค่า \Rsq{} ได้ไม่เกิน 0.045
การติดตามข่าวเชิงลบสัมพันธ์กับความวิตกกังวลมากกว่ามิติอื่น แต่ไม่มีนัยสำคัญ
ตัวแปรที่แยกคะแนนภาวะซึมเศร้าได้ชัดที่สุดคือความเพียงพอทางการเงิน (\etasq{}~= 0.17)
รองลงมาคือคุณภาพการนอน ผลนี้ชี้ว่าการดูแลสุขภาพจิตนักศึกษา
ควรให้ความสำคัญกับเรื่องค่าใช้จ่ายและการนอนมากกว่าเวลาที่ใช้หน้าจอ

\noindent\textbf{\textit{คำสำคัญ -{}-}} โซเชียลมีเดีย, สุขภาวะทางจิต, DASS-21, นักศึกษา,
ความเพียงพอทางการเงิน

\vspace{\baselineskip}
\noindent{\centering\bfseries ABSTRACT\par}

This study examined whether daily social media use, the type of content viewed,
and the following of negative news are related to depression, anxiety, and stress,
measured with the Thai version of the DASS-21.
Data were collected through an online survey of undergraduate students at Khon Kaen University.
Of 121 responses, 107 respondents who passed an attention check were analyzed.
Pearson correlation and hierarchical multiple regression showed that daily hours of social media use
were not significantly related to any dimension (\stat{r} from −0.15 to −0.09).
After controlling for sleep, academic pressure, grade point average, and financial sufficiency,
the social media variables added at most 0.045 to \Rsq.
Following negative news was related to anxiety more than to the other dimensions,
but not significantly.
Financial sufficiency separated depression scores most clearly (\etasq{}~= 0.17),
followed by sleep quality.
These results suggest that student mental health support should focus on money and sleep
rather than on screen time.

\noindent\textbf{Keywords -{}-} \textit{social media, mental well-being, DASS-21,
university students, financial sufficiency}
